\problem{Tridiagonal solver} \label{problem:tridiag}
    In this task we want to practice the principle of 
    Gaussian elimination on a simplified setup. Consider
    the invertible tridiagonal matrix
    \begin{equation}
        \mat{A} = \left( \begin{array}{cccccc}
            d_0 & u_0 & 0 & \dots & 0 & 0 \\
            l_1 & d_1 & u_1 & 0 & \dots &  \\
            0 & l_2 & d_2 & u_2 & 0 & \dots \\
            0 & 0 & \ddots & \ddots & \ddots & 0 \\
            0 & \dots &  & l_{N-2} & d_{N-2} & u_{N-2} \\
            0 & 0 & \dots & 0 & l_{N-1} & d_{N-1}
        \end{array} \right)
    \end{equation}
    (by convention $l_0 = 0, u_{N-1} = 0$) where we want to solve
    \begin{equation}
        \mat{A} \vec{x} = \vec{b}
    \end{equation}
    for $\vec{x}$.
    \begin{enumerate}
        \item Write the system in the form $(\mat{A}|\vec{b})$. Find the row operations which eliminate the lower diagonal, starting
        with $l_1$. \textit{Hint}: You start by subtracting a multiple of the first row from the second such that $l_1 = 0$ and continue downward.
        \item Once the lower diagonal is eliminated, the system can be solved bottom-up. Given the transformed system
        with eliminated lower diagonal, find an expression for the solution vector. \textit{Hint}: Start by finding an expression for 
        the last entry $x_{N-1}$ in $\vec{x} = (x_0, x_1, ..., x_{N-2}, x_{N-1})^T$, continue with $x_{N-2}$ and so on.
        \item Implement the tridiagonal solver in Python. The signature of the function as well as the 
        implementation of a matrix-vector product are given in Code-Snippets \ref{code:tridiag_task1} to \ref{code:tridiag_task3}. What is the
        scaling of your algorithm? Compare to $\mathcal{O}(N^3)$ of standard LU factorization.
    \end{enumerate}

    \begin{codebox}
        \begin{minted}{python3}
            """
            The tridiagonal matrix is given by
            [u_0, u_{N-2}], [d_0, d_{N-1}], [l_1, l{N-1}].
            By convention all arrays are of size N, with
            u_{N-1} = 0, l_0 = 0. This way the
            (l_i, d_i, u_i) are the values 
            from the i-th row.
            """
            import numpy as np
            import matplotlib.pyplot as plt
            from scipy.linalg import lu_factor, lu_solve

            def tri_vector_product(l, d, u, x):
                """
                Compute Ax where A is a tridiagonal
                matrix defined by l, d, u.
                """
                return x * d + np.roll(x, 1) * l + np.roll(x, -1) * u

            def tri_to_matrix(l, d, u):
                """
                Convert a tridiagonal matrix to a
                full matrix.
                """
                N = d.shape[0]
                return np.diag(d) + np.diag(l[1:], -1) + np.diag(u[:-1], 1)
        \end{minted}
        \caption{Python code for a matrix-vector product $\mat{A}\vec{x}$ where $\mat{A}$ is tridiagonal
        and signature for a tridiagonal solve of $\mat{A}\vec{x} = \vec{b}$. Continued in Code-Snippet \ref{code:tridiag_task2}.}
        \label{code:tridiag_task1}
    \end{codebox}

    \begin{codebox}
        \begin{minted}{python3}
            def tri_solve(l, d, u, b):
                """
                Solve Ax = b where A is a tridiagonal
                matrix defined by l, d, u.

                Args:
                    l, d, u: arrays of size N
                    b: array of size N

                Returns:
                    Solution to Ax = b of size N
                """
                N = diag.shape[0]

                # backward elimination of 
                # the lower diagonal
                for i in range(0, N - 1):
                    # TODO

                # forward substitution
                x = np.zeros_like(diag)
                x[-1] = # TODO

                for i in range(N - 2, -1, -1):
                    # TODO
                
                return x
        \end{minted}
        \caption{Continuation of Code-Snippet \ref{code:tridiag_task1}. Continued in Code-Snippet \ref{code:tridiag_task3}}
        \label{code:tridiag_task2}
    \end{codebox}

    \begin{codebox}
        \begin{minted}{python3}
            # test setup
            l = np.array([0.0, 2.0, 4.0, 6.0, 2.0, 6.0])
            d = np.array([7.0, 9.0, 10.0, 15.0, 9.0, 7.0])
            u = np.array([4.0, 6.0, 5.0, 8.0, 7.0, 0.0])

            # predefined solution
            x = np.array([5.0, 7.0, 8.0, 9.0, 5.0, 1.0])

            # tridiagonal representation to full matrix
            A = tri_to_matrix(l, d, u)

            # matrix vector product
            assert np.allclose(A @ x, tri_vector_product(l, d, u, x))
            b = tri_vector_product(l, d, u, x)

            # tridiagonal solve
            x_tilde = tri_solve(l, d, u, b)
            assert np.allclose(x, x_tilde)
        \end{minted}
        \caption{Continuation of Code-Snippet \ref{code:tridiag_task2}.}
        \label{code:tridiag_task3}
    \end{codebox}

\begin{solutionblock}
\part
    To eliminate the lower diagonal $l_i$, we proceed from $i=1$ to $N-1$. For each row $i$, we subtract a multiple of the modified row $i-1$ from row $i$. Let the modified diagonal elements be $d'_i$ and the modified right-hand side elements be $b'_i$.
    \begin{enumerate}
        \item For the first row ($i=0$), the elements remain unchanged:
            \begin{equation}
                d'_0 = d_0, \quad b'_0 = b_0
            \end{equation}
        \item For $i = 1, \dots, N-1$, we calculate the multiplier $m_i = \frac{l_i}{d'_{i-1}}$. The row operations are:
            \begin{itemize}
                \item $d'_i = d_i - m_i u_{i-1}$
                \item $b'_i = b_i - m_i b'_{i-1}$
            \end{itemize}
    \end{enumerate}
    The upper diagonal $u_i$ remains unchanged because the element above it in the 
    previous row's column is zero. After these operations, the matrix is upper bidiagonal.
    \part
    After the lower diagonal is eliminated, the system is in upper triangular form:
    \begin{equation}
        \left( \begin{array}{ccccc|c}
            d'_0 & u_0  &      &        &        & b'_0 \\
                & d'_1 & u_1  &        &        & b'_1 \\
                &      & \ddots & \ddots &        & \vdots \\
                &      &      & d'_{N-2} & u_{N-2} & b'_{N-2} \\
                &      &      &        & d'_{N-1} & b'_{N-1}
        \end{array} \right)
    \end{equation}
    We solve this using back-substitution starting from the last row:
    \begin{enumerate}
        \item For the last entry: $x_{N-1} = \frac{b'_{N-1}}{d'_{N-1}}$
        \item For $i = N-2, N-3, \dots, 0$:
        \begin{equation}
            d'_i x_i + u_i x_{i+1} = b'_i \implies x_i = \frac{b'_i - u_i x_{i+1}}{d'_i}
        \end{equation}
    \end{enumerate}
    \part
    The implementation of the Thomas algorithm (tridiagonal matrix solver) is given
    in Code-Snippet \ref{code:tridiag_sol}.

    \begin{codebox}
        \begin{minted}{python3}
            def tri_solve(l, d, u, b):
                N = d.shape[0]
                # Copy to avoid modifying original arrays
                d_copy = d.astype(float).copy()
                b_copy = b.astype(float).copy()

                # backward elimination of the lower diagonal (forward sweep)
                for i in range(0, N - 1):
                    m = l[i+1] / d_copy[i]
                    d_copy[i+1] -= m * u[i]
                    b_copy[i+1] -= m * b_copy[i]

                # forward substitution (back substitution sweep)
                x = np.zeros_like(d_copy)
                x[-1] = b_copy[-1] / d_copy[-1]

                for i in range(N - 2, -1, -1):
                    x[i] = (b_copy[i] - u[i] * x[i+1]) / d_copy[i]
                
                return x
        \end{minted}
        \caption{Thomas algorithm (tridiagonal matrix solver).}
        \label{code:tridiag_sol}
    \end{codebox}

\textbf{Scaling Analysis:}
The algorithm consists of two loops, each iterating $N$ times. Inside each loop, 
only a constant number of floating-point operations (additions, multiplications, divisions) 
are performed. Therefore, the complexity is $\mathcal{O}(N)$. 

Compared to the standard Gaussian elimination or LU factorization which 
scales as $\mathcal{O}(N^3)$, this solver is significantly more efficient 
for large tridiagonal systems, reducing both computation time and memory 
requirements (since we only store vectors of size $N$ rather 
than a matrix of size $N^2$).
\end{solutionblock}

\pagebreak